% Unit 2 guided notes -- 6 block lessons.
\documentclass{apchem-notes}

\apcunit{2}
\apcunittitle{Compound Structure and Properties}
\apcdoctitle{Guided Notes}
\apcsubtitle{CED 2.1--2.7 \textbullet\ Zumdahl \S8.1--8.13, \S9.1, \S10.4, \S10.7}

\begin{document}
\apctitleblock

% =====================================================================
\blocklesson{Bond Types and Potential Energy}{Wed Aug 26}

\begin{objectives}
  \item Classify bonding from electronegativity difference and element type.
  \item Read bond length and bond energy off a potential-energy curve.
\end{objectives}

\reading{Zumdahl \S8.1--8.3}{391--400}\quad\reading{\S8.8}{412--415}

\begin{warmup}[3]
  \item Configuration of \ce{Cl-}: \answerblank[3.5cm]{$[\ce{Ar}]$}
  \item Higher first IE: Mg or Al? \answerblank[6.7cm]{Mg (Al's 3p e$^-$ is
        higher-energy)}
  \item Formula: calcium + fluorine: \answerblank[2.5cm]{\ce{CaF2}}
\end{warmup}

\segment{INSTRUCTION A}{25}{A spectrum of bonds, not three boxes}

\notesection{Types of chemical bonds}{2.1}
\practice{4}

Bonding is a competition for electrons, scored by
\term{electronegativity} (EN):

\renewcommand{\arraystretch}{1.35}
\noindent\begin{tabular}{@{}p{3.4cm}p{3.6cm}p{6.2cm}@{}}
\toprule
\textbf{Bond type} & \textbf{Typical partners} & \textbf{Electron behavior}\\
\midrule
Nonpolar covalent & two identical nonmetals & shared
  \answerblank[2.2cm]{equally}\\
Polar covalent & different nonmetals & shared
  \answerblank[2.4cm]{unequally} --- partial charges $\delta^+/\delta^-$\\
Ionic & metal + nonmetal & \answerblank[2.6cm]{transferred} to the
  nonmetal\\
Metallic & metal + metal & delocalized ``\answerblank[2.6cm]{sea}'' over
  cations\\
\bottomrule
\end{tabular}

\begin{apcnote}
The AP exam treats these as a \emph{continuum}: greater $\Delta$EN $\to$
more ionic character. \ce{HF} is very polar covalent; \ce{NaCl} is mostly
ionic. Don't fight over borderline cases --- argue with $\Delta$EN.
\end{apcnote}

Rank the bond polarity, least to most: \ce{C-H}, \ce{C-F}, \ce{C-O}:
\answerblank[7.1cm]{\ce{C-H} $<$ \ce{C-O} $<$ \ce{C-F} ($\Delta$EN grows)}

\segment{GUIDED PRACTICE}{15}{Classify and justify}

For each substance: bond type + one-phrase justification.
\begin{enumerate}[leftmargin=*,itemsep=0.35em]
  \item \ce{Br2}: \answerblank[7.4cm]{nonpolar covalent --- identical atoms}
  \item \ce{K2O}: \answerblank[7.8cm]{ionic --- metal + nonmetal, large $\Delta$EN}
  \item \ce{NO}: \answerblank[8.9cm]{polar covalent --- two nonmetals, small $\Delta$EN}
  \item brass (Cu/Zn): \answerblank[6.5cm]{metallic --- delocalized electrons}
\end{enumerate}

\segment{INSTRUCTION B}{20}{The potential-energy curve}

\notesection{Why bonds form: the energy well}{2.2}
\practice{4}

\begin{center}
\begin{tikzpicture}
\begin{axis}[width=9.2cm, height=5.4cm,
    xlabel={internuclear distance (pm)}, ylabel={potential energy (kJ/mol)},
    xmin=30, xmax=300, ymin=-500, ymax=420,
    xtick={74}, xticklabels={74},
    ytick={0,-432}, yticklabels={0, $-432$},
    axis lines=left, every axis y label/.append style={rotate=0}]
  \addplot[seriesA, domain=48:295, samples=120]
    {432*((1-exp(-0.02*(x-74)))^2 - 1)};
  \draw[densely dotted, thick] (axis cs:74,-432) -- (axis cs:74,0);
  \draw[densely dotted, thick] (axis cs:30,-432) -- (axis cs:74,-432);
\end{axis}
\end{tikzpicture}
\end{center}

Reading the curve for \ce{H2}:
\begin{itemize}[leftmargin=*,itemsep=0.25em]
  \item At large distance, PE $\approx$ \answerblank[1.6cm]{0} --- atoms
        don't interact.
  \item As atoms approach, PE \answerblank[2.2cm]{decreases}: each nucleus
        attracts the other's electron density.
  \item The minimum sits at the \term{bond length}
        (\answerblank[1.8cm]{74 pm} for \ce{H2}); its depth is the
        \term{bond energy} (\answerblank[2.5cm]{432 kJ/mol}).
  \item At very short distance PE shoots up:
        \answerblank[5.3cm]{nucleus--nucleus repulsion} dominates.
\end{itemize}

\begin{apctrap}
The minimum is a \emph{balance} of attraction and repulsion --- not ``where
attraction stops.'' And breaking a bond \emph{requires} energy (climb out of
the well); forming one \emph{releases} it. Bond breaking is never
exothermic.
\end{apctrap}

\segment{APPLICATION}{20}{Compare two wells}

\ce{N2}'s well: minimum at 110 pm, depth 945 kJ/mol. \ce{H2}: 74 pm,
432 kJ/mol.

\begin{enumerate}[leftmargin=*,itemsep=0.4em]
  \item Which bond is stronger, and what feature of the curve says so?
        \answerlines[2]{\ce{N2} --- deeper well (945 vs 432 kJ/mol): more
        energy to climb out.}
  \item \ce{N2}'s bond is a triple bond. Connect that to both numbers.
        \answerlines[2]{Three shared pairs pull the nuclei closer would-be
        --- but N is larger than H, so compare like to like: vs a N--N
        single bond, the triple is much shorter and far deeper.}
  \item On one set of axes, sketch curves for \ce{F2} (bond energy
        155 kJ/mol, length 142 pm) and \ce{H2}. Label which is which.
        \workspace[3.6cm]
        \keyonly{\begin{apcanswer}\ce{F2}: shallower well (155), minimum
        farther right (142 pm). \ce{H2}: deeper (432), minimum at 74
        pm.\end{apcanswer}}
\end{enumerate}

\begin{exitticket}
On a PE curve, what two pieces of information does the minimum give you?
\answerlines[2]{Its x-position is the equilibrium bond length; its depth
below zero is the bond energy.}
\end{exitticket}

% =====================================================================
\blocklesson{Ionic Solids, Metals, and Alloys}{Fri Aug 28}

\begin{objectives}
  \item Rank lattice energies using Coulomb's law.
  \item Explain metallic properties from the electron-sea model; classify
        alloys.
\end{objectives}

\reading{Zumdahl \S8.5}{404--408}\quad\reading{\S10.4, 10.7}{505--511, 520--523}

\begin{warmup}[3]
  \item Most polar bond: \ce{H-F}, \ce{H-Cl}, or \ce{H-Br}?
        \answerblank[3cm]{\ce{H-F}}
  \item On a PE curve, bond energy is the \answerblank[3.5cm]{well depth}.
  \item Smaller radius: \ce{Na+} or \ce{Mg^2+}? Why (5 words)?
        \answerblank[6.6cm]{\ce{Mg^2+} --- more protons, same e$^-$}
\end{warmup}

\segment{INSTRUCTION A}{25}{Ionic solids and lattice energy}

\notesection{The lattice and its energy}{2.3}
\practice{6}

An ionic solid is a 3-D \term{lattice}: each cation surrounded by anions and
vice versa --- no ``\ce{NaCl} molecule,'' just a repeating ratio. The
\term{lattice energy} is the energy change when gaseous ions assemble into
the solid; its magnitude follows Coulomb's law:
\[
  |E| \propto \answerblank[3.5cm]{$\dfrac{|q_1 q_2|}{r}$}
\]

Decision order for ranking:
\begin{enumerate}[leftmargin=*,itemsep=0.2em]
  \item Compare \answerblank[3.4cm]{charge products} first --- they dominate.
  \item Tie? Compare \answerblank[3.2cm]{ionic radii} --- smaller ions,
        stronger lattice.
\end{enumerate}

\begin{workedexample}[Worked example: rank MgO, NaF, NaCl]
Charge products: MgO $(2\times2)=4$; NaF and NaCl $(1\times1)=1$. So MgO is
far strongest. Tiebreak NaF vs NaCl: \ce{F-} is smaller than \ce{Cl-}, so
NaF $>$ NaCl. Final: $\ce{MgO} \gg \ce{NaF} > \ce{NaCl}$.
\end{workedexample}

Consequences you can be asked to explain: high melting points
(strong \answerblank[4.9cm]{electrostatic attractions} throughout the
lattice), brittleness (shifting layers aligns
\answerblank[3.5cm]{like charges} $\to$ repulsion $\to$ cleave), and
conductivity only when \answerblank[3.8cm]{molten or dissolved} (ions must
be mobile).

\segment{GUIDED PRACTICE}{15}{Lattice ranking}

Rank by melting point, highest first, one reason each: \ce{KCl}, \ce{CaO},
\ce{KBr}.
\answerlines[3]{\ce{CaO} (charges $2\times2$) $>$ \ce{KCl} $>$ \ce{KBr}
(\ce{Cl-} smaller than \ce{Br-}). Charges first, then radius.}

\segment{INSTRUCTION B}{20}{Metals and alloys}

\notesection{The electron sea}{2.4}
\practice{1}

Metallic bonding: cations in fixed positions, valence electrons
\term{delocalized} over the whole sample. This one model explains:

\begin{itemize}[leftmargin=*,itemsep=0.25em]
  \item conductivity: \answerblank[6.0cm]{mobile electrons carry charge}
  \item malleability: layers slide without
        \answerblank[4.5cm]{aligning like charges} (contrast: ionic solids
        shatter)
  \item luster, thermal conduction: same mobile sea.
\end{itemize}

\notesub{Alloys}

\begin{itemize}[leftmargin=*,itemsep=0.25em]
  \item \term{Substitutional}: similar-size atoms \answerblank[3.2cm]{replace}
        host atoms (brass: Zn for Cu).
  \item \term{Interstitial}: small atoms fill \answerblank[2.6cm]{gaps}
        between hosts (steel: C in Fe). Interstitials pin sliding layers
        $\to$ \answerblank[4.5cm]{harder, less malleable}.
\end{itemize}

\segment{APPLICATION}{20}{Model-based predictions}

\begin{enumerate}[leftmargin=*,itemsep=0.4em]
  \item Solid \ce{NaCl} does not conduct; molten \ce{NaCl} does. Solid Cu
        conducts either way. Explain both with particle-level models.
        \answerlines[3]{In solid \ce{NaCl}, ions are locked in the lattice
        (no mobile charges); melting frees the ions. Cu's charge carriers
        are delocalized electrons, mobile in any phase.}
  \item Pure iron bends easily; steel (Fe + small \% C) resists. Why, at the
        particle level? \answerlines[2]{Carbon atoms sit in interstitial
        holes and block the Fe layers from sliding past one another.}
\end{enumerate}

\begin{exitticket}
Predict which has the higher melting point, \ce{LiF} or \ce{CsBr}, with the
two-step Coulombic argument. \answerlines[2]{Same charge product (1); Li$^+$
and F$^-$ are much smaller $\to$ shorter $r$ $\to$ stronger attraction $\to$
\ce{LiF}.}
\end{exitticket}

% =====================================================================
\blocklesson{Lewis Diagrams}{Mon Aug 31}

\begin{objectives}
  \item Draw Lewis diagrams for molecules and polyatomic ions quickly and
        correctly.
  \item Know the octet exceptions the AP exam actually tests.
\end{objectives}

\reading{Zumdahl \S8.9--8.11}{415--423}

\begin{warmup}[4]
  \item Stronger lattice: \ce{MgS} or \ce{NaCl}? Why?
        \answerblank[5.8cm]{\ce{MgS} --- $2\times2$ charge product}
  \item Valence electrons in \ce{P}: \answerblank[1.5cm]{5}
  \item Valence electrons in \ce{SO4^2-} total:
        \answerblank[3cm]{$6+4(6)+2 = 32$}
  \item Steel is which alloy type? \answerblank[3cm]{interstitial}
\end{warmup}

\segment{INSTRUCTION A}{25}{The five-step draw}

\notesection{Lewis diagrams}{2.5}
\practice{1}

\begin{enumerate}[leftmargin=*,itemsep=0.2em]
  \item Count \emph{total} valence electrons
        (\answerblank[7.1cm]{add for anions, subtract for cations}).
  \item Skeleton: \answerblank[6.8cm]{least electronegative atom central}
        (never H); connect with single bonds.
  \item Complete outer atoms' octets (H gets \answerblank[1.2cm]{2}).
  \item Leftovers go on the \answerblank[2.6cm]{central atom}.
  \item Central atom short? Convert lone pairs to
        \answerblank[3.4cm]{multiple bonds}.
\end{enumerate}

\begin{workedexample}[Worked example: \ce{CO2} in 20 seconds]
$4 + 2(6) = 16$ e$^-$. Skeleton O--C--O (12 left). Octets on the O's uses
all 16 --- C has only 4. Convert one lone pair from each O:
\chemfig{O=C=O} with two lone pairs on each oxygen --- every atom has an
octet, 16 e$^-$ total. Always end with the two-count check: electrons AND
octets.
\end{workedexample}

\notesub{The exceptions worth knowing}

\begin{itemize}[leftmargin=*,itemsep=0.25em]
  \item \term{Odd-electron} species (\ce{NO}, \ce{NO2}): someone keeps an
        unpaired electron.
  \item \term{Incomplete octet}: \ce{BF3} --- B holds \answerblank[1.4cm]{6}
        happily.
  \item \term{Expanded octet}: period-3+ central atoms (\ce{SF6},
        \ce{PCl5}) may hold more than 8 --- possible because $n\ge3$.
\end{itemize}

\segment{GUIDED PRACTICE}{15}{Draw four}

Draw Lewis diagrams (show electron count first):
\ce{NH3}, \ce{HCN}, \ce{ClO-}, \ce{PF5}
\workspace[4.6cm]
\keyonly{\begin{apcanswer}\ce{NH3}: 8 e$^-$, N central, 1 lone pair.
\ce{HCN}: 10 e$^-$, H-C$\equiv$N with lone pair on N. \ce{ClO-}: 14 e$^-$,
single bond, 3 lone pairs each, in brackets with $-$. \ce{PF5}: 40 e$^-$,
expanded octet, 5 single bonds, 3 lone pairs per F.\end{apcanswer}}

\segment{INSTRUCTION B}{20}{Bond order, length, strength}

\notesection{What the drawing predicts}{2.5}
\practice{4}

More shared pairs $\to$ higher \term{bond order} $\to$
\answerblank[2.6cm]{shorter} and \answerblank[2.6cm]{stronger} bond.

C--O bond in three species:
\renewcommand{\arraystretch}{1.3}
\noindent\begin{tabular}{@{}llll@{}}
\toprule
 & \ce{CH3OH} & \ce{CO2} & \ce{CO}\\
\midrule
bond order & \answerblank[1.4cm]{1} & \answerblank[1.4cm]{2} &
  \answerblank[1.4cm]{3}\\
length & longest & middle & shortest\\
\bottomrule
\end{tabular}

\segment{APPLICATION}{20}{Diagnose broken diagrams}

Each diagram below has one error; name it and fix it.
\begin{enumerate}[leftmargin=*,itemsep=0.4em]
  \item \ce{H2O} drawn with H--O--H and \emph{two} lone pairs on each H.
        \answerlines[2]{H can hold only 2 electrons --- lone pairs belong on
        O (2 pairs), giving 8 total.}
  \item \ce{CO3^2-} drawn with 22 electrons.
        \answerlines[2]{Count is $4+18+2 = 24$; the $2-$ charge ADDS two
        electrons.}
  \item \ce{NO+} drawn with a double bond and 12 electrons.
        \answerlines[2]{$5+6-1 = 10$ e$^-$ --- cation subtracts one; triple
        bond, one lone pair each.}
\end{enumerate}

\begin{exitticket}
Total valence electrons in \ce{PO4^3-}, and is an expanded octet allowed on
P? \answerlines[2]{$5+4(6)+3 = 32$; yes --- P is period 3 ($n=3$), though
the all-single-bond structure already satisfies octets.}
\end{exitticket}

% =====================================================================
\blocklesson{Formal Charge and Resonance}{Wed Sep 2}

\begin{objectives}
  \item Compute formal charge and use it to select the best structure.
  \item Draw complete resonance sets and state their physical meaning.
\end{objectives}

\reading{Zumdahl \S8.12}{423--427}

\begin{warmup}[3]
  \item Valence electrons in \ce{NO3-}: \answerblank[2.9cm]{$5+18+1=24$}
  \item Shortest C--O bond: \ce{CO}, \ce{CO2}, or \ce{CH3OH}?
        \answerblank[2.5cm]{\ce{CO}}
  \item Central atom in \ce{OCS} (carbon compound):
        \answerblank[5.1cm]{C (least electronegative)}
\end{warmup}

\segment{INSTRUCTION A}{25}{Formal charge: the tiebreaker}

\notesection{Formal charge}{2.6}
\practice{5}

\[
  FC = (\text{valence e}^-) -
  \answerblank[6.6cm]{$(\text{nonbonding e}^-) - \tfrac12(\text{bonding e}^-)$}
\]

Selection rules --- the best structure has:
\begin{enumerate}[leftmargin=*,itemsep=0.2em]
  \item formal charges closest to \answerblank[1.4cm]{zero};
  \item any negative FC on the \answerblank[5.5cm]{more electronegative atom};
  \item FCs summing to the species' \answerblank[3.2cm]{overall charge}
        (always --- this is a check, not a preference).
\end{enumerate}

\begin{workedexample}[Worked example: \ce{OCN-} skeleton choice]
Candidate skeletons: O--C--N vs C--O--N. For O--C--N with a C$\equiv$N
triple: FC(O)$=-1$, FC(C)$=0$, FC(N)$=0$ --- sum $-1$ \checkmark, minimal
charges, negative on electronegative O. The C--O--N skeletons force larger
charges (e.g.\ $+1$ on O). Carbon takes the center; formal charge proves it.
\end{workedexample}

\segment{GUIDED PRACTICE}{15}{Compute and choose}

For \ce{CO2} drawn two ways --- (i) two double bonds; (ii) single + triple:
\begin{enumerate}[leftmargin=*,itemsep=0.35em]
  \item FCs for (i): \answerblank[4.5cm]{0, 0, 0}
  \item FCs for (ii): \answerblank[5.2cm]{$-1$ (single-bond O), 0, $+1$}
  \item Which is better and why?
        \answerblank[6.5cm]{(i) --- all formal charges zero}
\end{enumerate}

\segment{INSTRUCTION B}{20}{Resonance}

\notesection{When one drawing isn't enough}{2.6}
\practice{1}

If the electron-pushing step could go more than one equivalent way, draw
\emph{all} of them connected by \answerblank[4.5cm]{double-headed arrows}.
The molecule is the \term{average} --- a resonance \term{hybrid} --- not a
movie flipping between forms.

\begin{workedexample}[Worked example: nitrate]
\ce{NO3-} (24 e$^-$): the double bond can sit on any of three oxygens ---
three equivalent forms. Experiment agrees: all three N--O bonds are
\emph{identical}, length between single and double, bond order $4/3$.
\end{workedexample}

\begin{apctrap}
``The bond switches back and forth'' scores zero. The measured structure is
one static average: equal bond lengths, fractional bond order, charge spread
over equivalent atoms. Cite the \emph{evidence} (equal lengths) when asked
what resonance means.
\end{apctrap}

\segment{APPLICATION}{20}{Full resonance analysis}

For carbonate, \ce{CO3^2-}:
\begin{enumerate}[leftmargin=*,itemsep=0.4em]
  \item Draw all three resonance forms. \workspace[3.6cm]
        \keyonly{\begin{apcanswer}Three forms, double bond rotating among
        the O's; 24 e$^-$ each; brackets with $2-$.\end{apcanswer}}
  \item Bond order of each C--O bond: \answerblank[2cm]{$4/3$}
  \item Predict how the C--O length compares with the C--O single bond in
        \ce{CH3OH} and the C=O in \ce{H2CO}.
        \answerlines[2]{Intermediate: shorter than the single bond in
        \ce{CH3OH}, longer than the true double bond in \ce{H2CO}.}
  \item Each O carries what average charge? \answerblank[2.5cm]{$-2/3$}
\end{enumerate}

\begin{exitticket}
\ce{NO2-} has two resonance forms. What does that predict about its two N--O
bonds, and what experiment would confirm it? \answerlines[2]{Equal length,
order $3/2$, between single and double; confirmed by measuring both bond
lengths (e.g.\ diffraction) and finding them identical.}
\end{exitticket}

% =====================================================================
\blocklesson{VSEPR: Shapes from Repulsion}{Fri Sep 4}

\begin{objectives}
  \item Predict electron-domain and molecular geometries through 6 domains.
  \item Explain bond-angle deviations with lone-pair repulsion.
\end{objectives}

\reading{Zumdahl \S8.13}{427--440}

\begin{warmup}[3]
  \item FC of S in \ce{SO4^2-} (all single bonds):
        \answerblank[2cm]{$+2$}
  \item Resonance forms of \ce{NO3-}: how many?
        \answerblank[1.5cm]{3}
  \item Bond order of N--O in \ce{NO3-}: \answerblank[2cm]{$4/3$}
\end{warmup}

\segment{INSTRUCTION A}{25}{Domains first, shape second}

\notesection{VSEPR}{2.7}
\practice{1}

Electron domains (bonding groups + lone pairs) around the central atom
repel and spread to maximum separation. A double bond counts as
\answerblank[3.2cm]{ONE domain}.

\renewcommand{\arraystretch}{1.3}
\noindent\begin{tabular}{@{}clll@{}}
\toprule
\textbf{Domains} & \textbf{Domain geometry} & \textbf{Angles} & \textbf{Example}\\
\midrule
2 & linear & 180$^\circ$ & \ce{CO2}\\
3 & trigonal planar & 120$^\circ$ & \ce{BF3}\\
4 & tetrahedral & 109.5$^\circ$ & \ce{CH4}\\
5 & trigonal bipyramidal & 90/120$^\circ$ & \ce{PCl5}\\
6 & octahedral & 90$^\circ$ & \ce{SF6}\\
\bottomrule
\end{tabular}

\notesub{Lone pairs bend the name}

Molecular geometry names only the \emph{atoms}:

\renewcommand{\arraystretch}{1.3}
\noindent\begin{tabular}{@{}llll@{}}
\toprule
\textbf{Domains (LP)} & \textbf{Molecular geometry} & \textbf{Angle} & \textbf{Example}\\
\midrule
4 (1 LP) & \answerblank[3.9cm]{trigonal pyramidal} & $\approx107^\circ$ & \ce{NH3}\\
4 (2 LP) & \answerblank[2.4cm]{bent} & $\approx104.5^\circ$ & \ce{H2O}\\
3 (1 LP) & bent & $<120^\circ$ & \ce{SO2}\\
\bottomrule
\end{tabular}

Angles shrink because lone pairs repel \answerblank[3.4cm]{more strongly}
than bonding pairs (held closer to one nucleus, fatter cloud).

\segment{GUIDED PRACTICE}{15}{Shape table sprint}

Lewis $\to$ domains $\to$ geometry $\to$ angle:
\begin{enumerate}[leftmargin=*,itemsep=0.35em]
  \item \ce{PF3}: \answerblank[9.4cm]{4 domains (1 LP) --- trigonal pyramidal, $\approx107^\circ$}
  \item \ce{CO3^2-}: \answerblank[7cm]{3 domains --- trigonal planar, 120$^\circ$}
  \item \ce{HCN}: \answerblank[7cm]{2 domains --- linear, 180$^\circ$}
  \item \ce{SF6}: \answerblank[7cm]{6 domains --- octahedral, 90$^\circ$}
\end{enumerate}

\segment{INSTRUCTION B}{20}{Explaining angles like a grader}

\notesection{The three-sentence angle answer}{2.7}
\practice{6}

Full-credit template, using \ce{H2O} vs \ce{CH4}:
\begin{enumerate}[leftmargin=*,itemsep=0.2em]
  \item Count: both have \answerblank[2.6cm]{4 domains} (tetrahedral
        arrangement).
  \item Distinguish: water has \answerblank[2.6cm]{2 lone pairs}; methane
        none.
  \item Conclude: lone pairs compress the H--O--H angle from 109.5$^\circ$
        to $\approx$\answerblank[2cm]{104.5$^\circ$}.
\end{enumerate}

\begin{apctrap}
Naming the shape is 1 point; the \emph{angle comparison} earns the rest.
``Bent, about 104.5$^\circ$, because two lone pairs repel more strongly than
bonding pairs'' --- practice the whole sentence.
\end{apctrap}

\segment{APPLICATION}{20}{Shapes in contrast}

\begin{enumerate}[leftmargin=*,itemsep=0.4em]
  \item \ce{CO2} is linear; \ce{SO2} is bent. Same atom count --- explain
        the difference completely. \answerlines[3]{C has 2 domains (two
        double bonds, no lone pairs) $\to$ 180$^\circ$. S has 3 domains
        (double bond, single-ish bond, lone pair) $\to$ bent,
        $<120^\circ$ --- the lone pair pushes the O's together.}
  \item Predict the geometry and angle of \ce{NO3-} and reconcile with its
        resonance. \answerlines[2]{Trigonal planar, 120$^\circ$ exactly ---
        resonance makes all three bonds identical, so no distortion.}
\end{enumerate}

\begin{exitticket}
\ce{NH3} vs \ce{BF3}: both ``\ce{XY3}.'' Different shapes --- why?
\answerlines[2]{N has 4 domains (1 lone pair) $\to$ trigonal pyramidal
$\approx107^\circ$; B has only 3 domains (6 e$^-$, no lone pair) $\to$
trigonal planar 120$^\circ$.}
\end{exitticket}

% =====================================================================
\blocklesson{Hybridization and Molecular Polarity}{Mon Sep 7 (confirm --- Labor Day)}

\begin{objectives}
  \item Assign sp, sp$^2$, sp$^3$ hybridization from electron domains; count
        $\sigma$ and $\pi$ bonds.
  \item Decide molecular polarity from shape + bond dipoles.
\end{objectives}

\reading{Zumdahl \S9.1}{455--466}

\begin{warmup}[4]
  \item Geometry of \ce{SO3}: \answerblank[3.5cm]{trigonal planar}
  \item Angle in \ce{H2S} vs 109.5$^\circ$: \answerblank[4.4cm]{smaller
        (2 lone pairs)}
  \item Bond order in \ce{O3}: \answerblank[2cm]{$3/2$}
  \item Lattice energy: bigger factor, charge or radius?
        \answerblank[2.5cm]{charge}
\end{warmup}

\segment{INSTRUCTION A}{25}{Hybridization: count domains, done}

\notesection{sp, sp$^2$, sp$^3$}{2.7}
\practice{1}

Hybridization is bookkeeping for domain count:

\renewcommand{\arraystretch}{1.3}
\noindent\begin{tabular}{@{}cll@{}}
\toprule
\textbf{Domains} & \textbf{Hybridization} & \textbf{Leftover p orbitals for $\pi$}\\
\midrule
2 & \answerblank[1.8cm]{sp} & 2\\
3 & \answerblank[1.8cm]{sp$^2$} & 1\\
4 & \answerblank[1.8cm]{sp$^3$} & 0\\
\bottomrule
\end{tabular}

\begin{apcnote}
The AP exam assesses ONLY sp, sp$^2$, sp$^3$. Five- and six-domain centers
(\ce{PCl5}, \ce{SF6}) still get shapes from VSEPR, but you will not be asked
their hybridization.
\end{apcnote}

\notesub{$\sigma$ and $\pi$}

Single bond = 1 $\sigma$. Double = $\sigma + \pi$. Triple =
$\sigma + 2\pi$. Count for \ce{H2C=CH2}:
\answerblank[3.4cm]{5 $\sigma$ + 1 $\pi$}.

\segment{GUIDED PRACTICE}{15}{Assign and count}

For \ce{HCN}: hybridization of C \answerblank[1.6cm]{sp}; total $\sigma$
\answerblank[1.2cm]{2}; total $\pi$ \answerblank[1.2cm]{2}.\\
For \ce{H2CO} (C central): C is \answerblank[1.6cm]{sp$^2$}; angles
$\approx$ \answerblank[1.6cm]{120$^\circ$}.\\
For \ce{CH3OH}: C is \answerblank[1.6cm]{sp$^3$}; O is
\answerblank[1.6cm]{sp$^3$} (2 lone pairs + 2 bonds).

\segment{INSTRUCTION B}{20}{Molecular polarity: the vector test}

\notesection{Polar molecule or not}{2.7}
\practice{6}

Two requirements, both necessary:
\begin{enumerate}[leftmargin=*,itemsep=0.2em]
  \item \answerblank[3.4cm]{polar bonds} exist ($\Delta$EN $\ne$ 0), AND
  \item the geometry does NOT \answerblank[3.2cm]{cancel} the bond dipoles.
\end{enumerate}

\begin{workedexample}[Worked example: \ce{CO2} vs \ce{H2O}]
\ce{CO2}: two strong C=O dipoles, but linear geometry points them exactly
opposite --- vector sum zero --- \textbf{nonpolar}. \ce{H2O}: two O--H
dipoles at 104.5$^\circ$ add to a net dipole toward O --- \textbf{polar}.
Same logic kills \ce{CCl4} (tetrahedral, cancels) and saves \ce{CHCl3}
(one different corner, doesn't cancel).
\end{workedexample}

\segment{APPLICATION}{20}{Unit synthesis --- structure to property}

For \ce{SO2} and \ce{SO3}:
\begin{enumerate}[leftmargin=*,itemsep=0.4em]
  \item Geometry of each: \answerblank[9.1cm]{\ce{SO2} bent ($<120^\circ$);
        \ce{SO3} trigonal planar (120$^\circ$)}
  \item One is polar, one is not --- assign and justify with the vector
        argument. \answerlines[2]{\ce{SO2} polar: bent shape leaves a net
        dipole. \ce{SO3} nonpolar: three equal dipoles at 120$^\circ$
        cancel.}
  \item Hybridization of S in each: \answerblank[4.4cm]{both sp$^2$ (3
        domains)}
\end{enumerate}

\begin{exitticket}
\ce{NF3} and \ce{BF3}: which is polar? Full argument --- shape, dipoles,
cancellation. \answerlines[3]{\ce{NF3}: trigonal pyramidal (lone pair on N),
three N--F dipoles don't cancel $\to$ polar. \ce{BF3}: trigonal planar,
symmetric, cancels $\to$ nonpolar.}
\end{exitticket}

\end{document}
